\begin{frame}{자주 묻는 질문 (2/2)}
  \textbf{Q. 라벨을 왜 굳이 $-1/+1$로 바꾸나요?}
  \begin{itemize}
    \item $y^{(i)}(w^\top x^{(i)}+b)$ \textbf{하나의 식}으로
      ``올바른 클래스 쪽에 있는가''와 ``얼마나 여유 있게 있는가''를
      동시에 표현하기 위한 \textbf{표기상의 편의}.
    \item $y=\pm1$이면 ``곱이 클수록 좋다''는 하나의 부등식으로 통일 —
      $0/1$ 라벨이었다면 클래스마다 부등호 방향이 다른 두 식을
      따로 다뤄야 했다.
  \end{itemize}
  \vskip0.5em
  \textbf{Q. 로지스틱회귀와 같은 $w^\top x+b=0$인데, 뭐가 다른가요?}
  \begin{itemize}
    \item \textbf{어떤} 경계를 뽑을지가 다르다:
      로지스틱회귀는 \textbf{모든 점}이 손실에 기여해 분포 전체에 민감
      (매끄러운 경계), SVM은 마진 바깥 점의 손실이 정확히 0이라
      오직 서포트 벡터만 경계를 결정(각진 경계).
    \item \textbf{이상치 견고성}: 이상치가 서포트 벡터가 안 되는 한
      SVM은 무시; 로지스틱회귀는 아무리 멀어도 아주 작은 기울기로
      손실에 영향 — 서로 다른 트레이드오프.
  \end{itemize}
\end{frame}
